\documentclass[ch3.tex]{subfiles}
%subfile : ch3_2.tex
\begin{document}
	\section{Generalization of Submetacompactness}  
	\label{sec:3-2}


{\definition{} \label{def:3-2-1}
	\textup{(Smith [1975]\textsuperscript{\cite{Smith1975}})} 
An open cover $\mathscr{U}$ of $X$ is a 
\uwave{weak $\overline{\theta}$-cover}
\index{weak $\overline{\theta}$-cover| \S\ref{sec:3-2}\hfill}
if $\mathscr{U}$ has an open refinement $\mathscr{G} = \bigcup_{n \in \omega} \mathscr{G}_n$ where each $\mathscr{G}_n = \{G_{n, \alpha} : \alpha \in \Gamma_n\}$ such that

\textup{(i)}
For each $x \in X$, there exists $m \in \omega$ such that 
$0 < |\operatorname{ind}(\mathscr{G}_m)_x| < \omega$.

\textup{(ii)}
$\mathscr{G}^* = \{\bigcup \mathscr{G}_n : n \in \omega\}$ is a point-finite open cover of $X$.

\noindent The family $\mathscr{G}$ is called a 
\uwave{weak $\overline{\theta}$-refinement}
\index{weak $\overline{\theta}$-refinement| \S\ref{sec:3-2}\hfil}
of $\mathscr{U}$.  $X$ is 
\uwave{$\lambda$-weakly $\overline{\theta}$-refinable}
\index{$\lambda$-weakly $\overline{\theta}$-refinable| \S\ref{sec:3-2}\hfil}
if every open $\lambda$-cover is a weak $\overline{\theta}$-cover.
}


{\definition{} \label{def:3-2-2}
\textup{(i)}
An open cover $\mathscr{U}$ of $X$ is a $D\omega$-cover if $\mathscr{U}$ 
has a refinement $\mathscr{F} = \bigcup_{n \in \omega} \mathscr{F}_n$ where each 
$\mathscr{F}_n$ is a family of relatively discrete and closed in the subspace 
$X \setminus\bigcup_{i< m}(\cup \mathscr{F_i})$, and for any $m\ne n$,
$(\cup\mathscr{F}_m)\bigcap(\cup\mathscr{F}_n) = \emptyset$.
The family $\mathscr{F}$ is called a 
\uwave{$D\omega$-refinement}
\index{$D\omega$-refinement| \S\ref{sec:3-2}\hfill}
of $\mathscr{U}.$

\textup{(ii)}
\textup{(Liu [1977]\textsuperscript{\cite{Liu}})} 
$X$ is 
\uwave{$\lambda$-strictly quasi-paracompact}
\index{$\lambda$-strictly quasi-paracompact| \S\ref{sec:3-2}\hfill}
or 
\uwave{$D\omega$-refinable}
\index{$D\omega$-refinable| \S\ref{sec:3-2}\hfill}
if every open $\lambda$-cover of $X$ is a $D\omega$-cover.
}

{\definition{}  \label{def:3-2-3}
\textup{(Smith [1980]\textsuperscript{\cite{Smith1980})}} 
Let $\omega^2 = \omega\cdot\omega$ be an ordinal product, $\xi\in\{\omega, \omega^2\}$. 
The properties P will be discrete \textup{(D)}, locally finite and closed \textup{(C)},.

\textup{(i)} An open cover $\mathscr{U}$ of$X$ is a 
\uwave{$B(P, \xi)$-cover}
\index{$B(P, \xi)$-cover| \S\ref{sec:3-2}\hfill}
if $\mathscr{U}$ has a refinement $\mathscr{E} = \bigcup_{\beta < \xi} \mathscr{E}_\beta$, which satisfies

}

Here is the transcription of the handwritten notes from page 33:

---

(a). $(\bigcup\mathcal{E}_\beta) \cap (\bigcup\mathcal{E}_{\beta'}) = \Phi \quad \forall \beta \ne \beta'$;

(b) for each $\beta < \xi$, $\mathcal{E}_\beta$ is a relatively LF family of closed subsets of the subspace $X - \bigcup_{\mu<\beta}(\bigcup\mathcal{E}_\mu)$;

(c). for each $\beta < \xi$, $\bigcup_{\mu<\beta}(\bigcup\mathcal{E}_\mu)$ is closed in $X$.

For the case $\mathcal{P} = \mathcal{C}$, we require $\mathcal{E}$ to be a one-to-one partial refinement of $\mathcal{U}$.

The family $\mathcal{E}$ is called a $B(\mathcal{P}, \xi)$-refinement of $\mathcal{U}$.

(ii) $X$ is $\mathfrak{m}$-$B(\mathcal{P}, \xi)$-refinable if every open $\mathfrak{m}$-cover of $X$ has a $B(\mathcal{P}, \xi)$-refinement.

It is easy to see that $X$ is $\mathfrak{m}$-strictly quasi-paracompact iff $X$ is $\mathfrak{m}$-$B(\mathcal{D}, \omega)$-refinable.

---

3.2.4. Lemma (Liu 1977).

A point-finite open cover of $X$ is a $D\omega$-cover.

**Proof.** Suppose $\mathcal{U} = \{ U_\alpha : \alpha \in A \}$ is a point-finite open cover of $X$. For $n < \omega$, let

$$D_n = \{ x \in X : \text{ord}(x, \mathcal{U}) \le n \}, \quad D_0 = \Phi$$

Then $D_n$ is closed. Let $F(\phi, 0) = \Phi$. For $n \ge 1$ and $s \in [A]^n$, let

$$F(s, n) = D_n \setminus \bigcup_{\alpha \in A \setminus s} U_\alpha$$

$$\mathcal{F}_n = \{ F(s, n) : s \in [A]^n \}$$

We show that $\mathcal{F} = \bigcup_{n<\omega}\mathcal{F}_n$ is a $D\omega$-refinement of $\mathcal{U}$. Let $n < \omega$, $Y = X - \bigcup_{i<n}(\bigcup\mathcal{F}_i)$. To show that $\mathcal{F}_n\vert{}_Y$ is a discrete closed family in $Y$. By Lemma 1.1.4, we need only to show that

Here is the complete transcription of the handwritten notes from page 34 (marked as 33 at the bottom):

---

$\mathcal{F}_n\vert{}_Y$ is closed in $Y$, and

(2) $\mathcal{F}_n\vert{}_Y$ is discrete at each $x \in Y$.

To see (1), let $G = Y - \bigcup(\mathcal{F}_n\vert{}_Y)$. Suppose $\text{ord}(x, \mathcal{U}) = j$. Then $x \in D_j$. Pick $s \in [A]^j$ such that $x \in \bigcap_{\alpha \in s} U_\alpha = F(s, j) \subset \bigcup \mathcal{F}_j$.

Since $x \in G$, $j > n$. Then $Y - D_n$ is a neighborhood of $x$ in $Y$ such that $Y - D_n \subset G$. Hence $G$ is open in $Y$.

To see (2), let $x \in \bigcup(\mathcal{F}_n\vert{}_Y)$. Suppose $s_0 \in [A]^n$ such that $x \in F(s_0, n)$. Then $W = Y \cap \bigcap_{\alpha \in s_0} U_\alpha$ is a neighborhood of $x$ in $Y$. For each $s \in [A]^n \setminus \{s_0\}$, $W \cap F(s, n) = \Phi$. So $\mathcal{F}_n\vert{}_Y$ is discrete at $x$. $\square$

---

 3.2.5. Theorem. (Zhu 1984, Long 1986a)

A $\mathfrak{m}$-submetacompact space is $\mathfrak{m}$-strictly quasi-paracompact.

**Proof.** Assume $X$ is $\mathfrak{m}$-submetacompact and let $\mathcal{U} = \{ U_\alpha : \alpha < \kappa \}$ be an open cover of $X$. By Lemma 3.2.1, there is a countable closed cover $\{ F_n \}$ of $X$ such that each $\mathcal{U}\vert{}_{F_n}$ has a point-finite open refinement $\mathcal{U}_n$ in $F_n$. By Lemma 3.2.4, let $\mathcal{E}_n = \bigcup_{k<\omega} \mathcal{E}_{n,k}$ be a $D\omega$-refinement of $\mathcal{U}_n$ in $F_n$. For each $n, k$, $\mathcal{E}_{n,k}$ is a discrete closed family in $F_n - \bigcup_{i<k}(\bigcup \mathcal{E}_{n,i})$. Then each $\mathcal{E}_{n,k}$ is a discrete closed family in $X - \bigcup_{i<k}(\bigcup \mathcal{E}_{n,i})$.

Let $\varphi$ be a one-to-one function from $\omega \times \omega$ onto $\omega$ defined by $\varphi(n, k) = \frac{1}{2}[(n+k)^2 + 3n + k]$. Define $\mathcal{F}_m = \mathcal{E}_{n,k}$ if $m = \varphi(n, k)$. Then $\mathcal{F} = \bigcup_{m<\omega} \mathcal{F}_m$ is a $D\omega$-refinement of $\mathcal{U}$. $\square$

Here is the complete transcription of the handwritten notes from page 34 (marked as 33/34 at the bottom):

---

 3.2.6. Theorem. (Smith 1980)

Every $\mathfrak{m}$-$D\omega$-refinable space is $\mathfrak{m}$-weakly $\bar{\theta}$-refinable.

**Proof.** Suppose $\mathcal{U} = [U_\alpha : \alpha < \lambda]$ be an open cover of a $\mathfrak{m}$-$D\omega$-refinable space $X$. 
Let $\mathcal{F} = \bigcup_{n<\omega}\mathcal{F}_n$ be a $D\omega$-refinement of $\mathcal{U}$, where each $\mathcal{F}_n = \{F_{n,\alpha} : \alpha < \lambda\}$ is a discrete family of closed subsets of $X - \bigcup_{i m}$. (3) holds.
	
	It follows from (1)–(3) that $\mathcal{G}$ is a weakly $\bar{\theta}$-refinement of $\mathcal{U}$. $\square$

\end{document}
